% page 35 du fac-simile (image p-034 du scan)
% bloc 170.2 x 233.3 mm ; marge gauche 31.8 mm
\pgc{10.160mm}{0.000mm}{
\l{\cel{55}27}
\l{}
\l{}
\l{\sou{anapesto}. Verso en qua la silabi alternas yene : 2 kurta,}
\l{\cel{3}l longa. - DEFIRS.}
\l{}
\l{\sou{anarkio}. Stando sociala qua karakterizesas da la ne-existo di}
\l{\cel{3}guvernanto. - DEFIRS.}
\l{}
\l{\sou{anastigmata}. (\sou{Opt}.) Vitro, od okulo, sen astigmateso.- DEFIRS.}
\l{}
\l{\sou{anastomoso}.\cel{2}(\sou{Anat}.)\cel{2}Junturo di du vaskuli. - DEFIRS.}
\l{}
\l{\sou{anatemar}.\cel{2}(\sou{trans}.) (\sou{religio kristana})\cel{2}Frapar per la malediko}
\l{\cel{3}per qua la Eklezio katolika forpulsas ulu de sua medio.- DEFIRS.}
\l{}
\l{\sou{anatomio}. Cienco qua traktas la strukturo di la organi di nia}
\l{\cel{3}korpo per dissekar li. - DEFIRS.}
\l{}
\l{\sou{ancestro}. Un persono ek la serio de patri, avi, preavi, e c.,}
\l{\cel{3}retroire til epoko nedefinita. - EF.}
\l{}
\l{\sou{anchovo}\cel{2}(\sou{zool}.) Mikra maro-fisho qua, konservate en salo,}
\l{\cel{3}manjesas kom acesorajo repastala. - L. engraulis encrasicholus.}
\l{\cel{3}- DEFIRS.}
\l{}
\l{\sou{anciena}. Qua esis tala ja multa tempo ante la instanto konsiderata.}
\l{\cel{3}- DEFIS.}
\l{}
\l{\sou{andrinoplo}. Sorto di stofo kotona, tintita rede. - FIS.}
\l{}
\l{\sou{androceo}. (\sou{bot.}) Ensemblo de la stamini, de la organi maskula}
\l{\cel{3}di floro. - DEF.}
\l{}
\l{\sou{androdioika}. (\sou{bot.}) Qualeso di planto che qua la flori esas :}
\l{\cel{3}ici, maskulo - iti, hermafrodita, la du sorti esante sur}
\l{\cel{3}stipiti diversa. - DEFIS.}
\l{}
\l{\sou{androgina}. (\sou{bot.})\cel{2}Qua posedas e la organi maskulala e la}
\l{\cel{3}organi feminala, ma en du flori separita. - DEFIRS.}
\l{}
\l{\sou{andromonoika}.- (\sou{bot.)}\cel{3}Qualeso di planto che qua la flori maskula}
\l{\cel{3}e la flori hermafrodita esas sur la sama stipito.- DEFIS.}
\l{}
\l{\sou{anekdoto}. Kurta naracuro pri mikra eventajo savinda. - DEFIRS.}
\l{}
\l{\sou{anemio.}\cel{1}(\sou{patol}.) Deperiso produktata da la povresko di sango,}
\l{\cel{3}- DEFIRS.}
\l{}
\l{\sou{anemofila}. (\sou{biol.})\cel{2}Dicesas pri la planti che qui la dissemo di}
\l{\cel{3}poleno eventas per vento. - DEFIS.}
\l{}
\l{\sou{anemometro}.\cel{2}Instrumento per qua onu povas mezurar la rapideso-}
\l{\cel{3}grado di vento. - DEFIRS.}
}
